% Generated draft: requires mathematical and source review.
\documentclass[11pt]{article}
\usepackage[margin=25mm]{geometry}
\usepackage{fontspec}
\setmainfont{Latin Modern Roman}
\usepackage{amsmath,amssymb,mathtools,unicode-math}
\setmathfont{Latin Modern Math}
\usepackage{enumitem,needspace,xurl,hyperref}
\setlength{\parindent}{0pt}
\setlength{\parskip}{5pt}
\setlength{\emergencystretch}{3em}
\newcommand{\N}{\mathbb{N}}
\newcommand{\Z}{\mathbb{Z}}
\newcommand{\Q}{\mathbb{Q}}
\newcommand{\R}{\mathbb{R}}
\newcommand{\C}{\mathbb{C}}
\newcommand{\F}{\mathbb{F}}
\newcommand{\A}{\mathbb{A}}
\newcommand{\PP}{\mathbb{P}}
\newcommand{\E}{\mathbb{E}}
\newcommand{\eps}{\varepsilon}
\newcommand{\dd}{\,\mathrm{d}}
\begin{document}
\section*{UnsolvedMath ID: 30005256}
\textbf{Source problem code:} catalogue entry
\subsection{Definitions and mathematical statement}
The supplied record concerns a surface \(S\), its Teichm\"uller space \(\mathcal T(S)\), its mapping class group \(\operatorname{Mod}(S)\), and a Lie group \(G_0\) of rank \(\ell\), with exponents denoted by \((1,m_2,\ldots,m_\ell)\). The topology of \(S\), boundary conditions, class of groups \(G_0\), and precise conventions for these spaces are not supplied. These choices must be fixed before a fully quantified problem can be stated.

For a complex structure \(J\in\mathcal T(S)\), write \(\Sigma_J=(S,J)\), let \(K_J\) be its holomorphic cotangent line bundle, and let \(H^0(\Sigma_J,K_J^{\otimes d})\) denote the space of holomorphic sections of its \(d\)-th tensor power. The displayed formula in the source fragment is naturally read as
\[
E_J=\bigoplus_{i=2}^{\ell}H^0\!\left(\Sigma_J,K_J^{\otimes(m_i+1)}\right),
\qquad E\longrightarrow\mathcal T(S).
\]
The fragment describes \(E\) as a holomorphic vector bundle. The direct-sum notation above reconstructs a damaged transcription and requires checking against the cited report. The notation \(\operatorname{Hit}(S)\) refers to the Hitchin component under discussion; its dependence on the target group is suppressed in the supplied fragment. The record also requests a formulation for maximal components, but does not specify their target groups or the invariant and normalization defining maximality.

\textbf{Requested construction, with unresolved scope.} For each Hitchin or maximal component \(\mathcal C\) intended by the reconstructed statement, construct an independently described geometric parameter space \(P_{\mathcal C}\), with an action of \(\operatorname{Mod}(S)\), and a parametrization
\[
\Phi_{\mathcal C}:P_{\mathcal C}\longrightarrow\mathcal C
\]
compatible with changes of marking. If this compatibility means equivariance, its precise condition is
\[
\Phi_{\mathcal C}(\gamma\cdot p)
 =\gamma\cdot\Phi_{\mathcal C}(p)
\quad\text{for every }\gamma\in\operatorname{Mod}(S)
\text{ and }p\in P_{\mathcal C}.
\]
Here a set-theoretic parametrization would require bijectivity; the source record does not specify whether a homeomorphism, diffeomorphism, or another regularity condition is intended. Nor does it specify whether ``invariant'' means this equivariance or a construction on mapping-class-group quotients. If an individual component is not preserved by the full group, the acting subgroup or invariant union of components must also be specified.

The intended construction should address the failure of uniqueness of the equivariant conformal maps used in the parametrization attributed to Labourie. The record does not supply that construction, its map from \(E\), or the conjecture referenced by the original fragment. Consequently, neither \(E\) itself nor any replacement can be asserted here to provide the requested parametrization. The admissible geometric description of \(P_{\mathcal C}\) must also be specified to make the construction request substantive.

\textbf{Supplied status qualification.} The registry labels the problem partially solved. Its dated assessment says that the construction is surjective, uniqueness holds in rank two, and later primary work gives nonuniqueness in rank at least three, so the fragment does not yield a general invariant parametrization. These are supplied literature assessments, with their group and representation hypotheses unspecified here, rather than independently verified mathematical assertions. The record's labels \texttt{corrected\_verified} and its statement-verification claim likewise report supplied curation.
\subsection{Short English statement}
Find a geometric parametrization of the intended higher-rank Hitchin and maximal components that respects the mapping class group and accommodates the reported failure of uniqueness in Labourie's construction. The supplied fragment defines a candidate bundle over Teichm\"uller space but omits the referenced conjecture, the parametrizing map, and essential surface and group hypotheses. The precise scope and meaning of parametrization therefore require human review.
\subsection{Sources}
\begin{itemize}
\item Local supplied registry record: UnsolvedMath permanent ID \(30005256\), title ``Mapping-Class-Invariant Parametrization of Hitchin Components,'' source problem code \texttt{catalogue entry}, in the Oberwolfach Reports Open Problems collection. Repository registry locations are \texttt{data/unsolvedmath/problems.json} and \texttt{catalog/problems.json}; neither was independently inspected for this draft. No individual proposer is supplied; the proposed year is recorded as 2022.
\item Supplied source citation: \emph{Arbeitsgemeinschaft: Higher Rank Teichm\"uller Theory} (2023), Oberwolfach Reports, pp.~2687--2740, DOI 10.4171/owr/2022/47. The extracted fragment carries page 2715. The literature list labels this ``Oberwolfach Report 47/2022.'' \url{https://doi.org/10.4171/owr/2022/47}. The supplied record says its reconstructed statement was checked against the official report; this draft has not independently verified that claim. The 2023 citation date and 2022 report designation are retained as supplied.
\item Supplied literature reference: ``Labourie, Annals of Mathematics.'' \url{https://doi.org/10.4007/annals.2017.185.1.1}. No article title or theorem-level attribution is supplied.
\item Supplied literature reference: ``Sagman--Smillie, Annals of Mathematics.'' \url{https://annals.math.princeton.edu/articles/22504}. No article title or precise hypotheses for the reported nonuniqueness result are supplied.
\item Literature assessment checked at \texttt{2026-08-22}: the supplied curation labels the problem partially solved and reports surjectivity, rank-two uniqueness, and higher-rank nonuniqueness. This is supplied dated curation derived from the OWR assessment, not independently verified literature review. The record expressly requests independent verification of time-sensitive status and cited claims.
\end{itemize}
\end{document}
